\chapter{Desarrollo matemático}

\section{Algunos resultados a tener en cuenta}

En los capítulos 8, 9, y 10 del libro de texto~\cite{Wegner2024}
se presentan algunas características no directamente esperables de
las distribuciones más importantes en Estadística cuando se trabaja
en espacios de alta dimensión.
Sin embargo, algunos fenómenos, como la concentración de medida o la
ortogonalidad de vectores, pueden resultar de utilidad en Análisis de
Datos.

\section{Proyecciones aleatorias}

En esta sección introducimos el concepto de proyecciones aleatorias.
Sea $\left(\Omega,\Sigma,\mathbb{P}\right)$ un espacio de
probabilidad, en esta sección se introducen matrices aleatorias del
tipo
\begin{align*}
	U\colon\Omega & \longrightarrow\mathbb{R}^{k\times d} \\
	\omega        & \longmapsto
	\begin{bmatrix}
		u_{11}\left(\omega\right) & \cdots & u_{1d}\left(\omega\right) \\
		\vdots                    & \ddots & \vdots                    \\
		u_{k1}\left(\omega\right) & \cdots & u_{kd}\left(\omega\right)
	\end{bmatrix},
\end{align*}
donde
\begin{math}
	\forall\left(i,j\right)\in
	\left\{1,\dotsc,k\right\}\times\left\{1,\dotsc,d\right\}:
	u_{ij}\sim\mathcal{N}\left(0,1\right)
\end{math}
son variables aleatorias normales unidimensionales e independientes.
Esta matriz aleatoria será denotada por
$U\sim\mathcal{N}\left(\symbf{0}_{k\times d},1\right)$.
Algunas propiedades que tiene la matriz aleatoria están basadas en el
comportamiento de sus filas como vectores normales multivariantes.
Si se denotan para cada fila $i=1,\dotsc,k$ de $U$ como
$U_{i}={\left[u_{i1},\dotsc,u_{id}\right]}^{T}$, entonces todas
ellas son vectores aleatorios independientes con distribución
$\mathcal{N}\left(\symbf{0}_{d},I_{d}\right)$
(ver Proposición~\ref{prop:independence}) y se pueden explotar los
teoremas 8.2, 8.3 y 8.4 del capítulo 8 del libro de
texto~\cite{Wegner2024} y, a su vez, los resultados del capítulo 10,
cuando $d$ sea elevado.
Por ejemplo, si usamos el Teorema 10.7 de ortogonalidad gaussiana, se
tiene que las filas de la matriz $U$ normalizadas son
``aproximadamente'' ortogonales unas a otras.
Entonces, normalizando la matriz aleatoria $U$, se obtendrá una
matriz aleatoria $V$ tal que $VV^{T}≈I_{k}$.
Dicha propiedad nos recuerda a la propiedad de ortogonalidad en las
matrices de proyección basadas en álgebra lineal.
Una de las aplicaciones más conocidas de las proyecciones es que
permiten transformar conjuntos de datos en $\mathbb{R}^{d}$ en el
subespacio de $\mathbb{R}^{k}$ generado por las $k$ filas de la
matriz ortogonal de proyección.
Se debe tener en cuenta que si $k<d$ entonces la proyección es una
transformación lineal, que cuenta con esta buena propiedad de
ortonormalidad que evita la ``duplicidad'' de información en los
datos proyectados, pero que no necesariamente conserva las distancias
entre puntos del espacio original $\mathbb{R}^{d}$.
Esta conservación de las distancias originales es una propiedad útil
en contextos de clasificación, por tanto supervisada como no
supervisada, por ejemplo, en métodos como los $m$-vecinos más
próximos.
El objetivo de este capítulo es probar que las proyecciones basadas
en matrices aleatorias del tipo
$U\sim\mathcal{N}\left(\symbf{0}_{k\times d},1\right)$, además que
ella tiene la característica de la ortonormalidad, de forma
aproximada, se mantienen también, de forma aproximada, las distancias
entre datos en el espacio original.

\begin{theorem}[Teorema de la Proyección aleatoria]\label{thm:projection}
	Sea $\left(\Omega,\Sigma,\mathbb{P}\right)$ un espacio de
	probabilidad y, para $k<d$ sea
	$U\colon\Omega\to\mathbb{R}^{k\times d}$ una matriz aleatoria con
	$U\sim\mathcal{N}\left(\symbf{0}_{k\times d},1\right)$.
	Para
	$\symbf{x}\in\mathbb{R}^{d}\setminus\left\{\symbf{0}_{d}\right\}$ y
	$0<\varepsilon\leq 1$ tenemos
	\begin{equation*}
		\mathbb{P}
		\left[
			\left|
			\left\|
			U\symbf{x}
			\right\|-
			\sqrt{k}
			\left\|
			\symbf{x}
			\right\|
			\right|\leq
			\varepsilon\sqrt{k}
			\left\|
			\symbf{x}
			\right\|
			\right]\geq
		1-2\exp
		\Big(-k\frac{\varepsilon^{2}}{16}\Big).
	\end{equation*}
\end{theorem}

\begin{proof}
	En primer lugar, se debe tener en cuenta que el producto matriz por
	vector de $U$ por $\symbf{x} \in \mathbb{R}^d$ da como resultado un
	vector en $\mathbb{R}^k$ que contiene los productos escalares entre
	cada vector fila de la matriz $U$, denotados por $U_{i}$, con el
	vector $\symbf{x}$.
	Es decir,
	\begin{equation*}
		U\symbf{x}=
		{\left[
			\left\langle
			U_{1},\symbf{x}
			\right\rangle,
			\dotsc,
			\left\langle
			U_{k},\symbf{x}
			\right\rangle
			\right]}^{\top}.
	\end{equation*}
	Ahora, fijado
	$\symbf{x}\in\mathbb{R}^{d}\setminus\left\{\symbf{0}_{d}\right\}$,
	considere la coordenada $i$-ésima del vector
	$U\dfrac{\symbf{x}}{\left\|\symbf{x}\right\|}$:
	\begin{equation*}
		{\left(
			U\frac{\symbf{x}}{\left\|\symbf{x}\right\|}
			\right)}_{i}=
		\left\langle
		U_{i},
		\frac{\symbf{x}}{\left\|\symbf{x}\right\|}
		\right\rangle=
		\sum_{j=1}^{d}
		u_{ij}
		\frac{x_{i}}{\left\|\symbf{x}\right\|}.
	\end{equation*}
	Como
	\begin{math}
		\forall\left(i,j\right)\in
		\left\{1,\dotsc,k\right\}\times\left\{1,\dotsc,d\right\}:
		u_{ij}\sim\mathcal{N}\left(0,1\right)
	\end{math}, se tiene que
	\begin{math}
		{\left(
			U\frac{\symbf{x}}{\left\|\symbf{x}\right\|}
			\right)}_{i}\sim
		\mathcal{N}\left(0,1\right)
	\end{math}
	(ver Proposición~\ref{prop:linearcombination}) luego, junto a la
	independencia de los $u_{i j}$, se tiene
	\begin{math}
		U\frac{\symbf{x}}{\left\|\symbf{x}\right\|}\sim
		\mathcal{N}\left(\symbf{0}_{d},I_{d}\right)
	\end{math}
	(ver Proposición~\ref{prop:independence}).
	Por tanto, usando el Teorema 10.4 del libro de texto (anillo
	Gaussiano) se obtiene
	\begin{equation*}
		\mathbb{P}
		\left[
			\left|
			\left\|U\symbf{x}\right\|-
			\sqrt{k}\left\|\symbf{x}\right\|
			\right|\leq
			\varepsilon\sqrt{k}
			\left\|\symbf{x}\right\|
			\right]=
		\mathbb{P}
		\left[
			\left|
			\big\|
			U\frac{\symbf{x}}{\left\|\symbf{x}\right\|}
			\big\|-
			\sqrt{k}
			\right|\leq\varepsilon\sqrt{k}
			\right] \geq
		1-2\exp\Big(-k\frac{\varepsilon^{2}}{16}\Big).
	\end{equation*}
	En esta acotación se ha aplicado dicho teorema con $d=k$ y
	$\varepsilon=\varepsilon \sqrt{k}$.
	Nótese que se cumple la hipótesis
	$0<\varepsilon\sqrt{k}\leq\sqrt{k}$ dado que $0<\varepsilon\leq1$.
\end{proof}
Puesto que $\sqrt{k}$ es fijo a lo largo de toda la prueba del
Teorema~\ref{thm:projection}, una reformulación más habitual de este
resultado es
\begin{equation}
	\mathbb{P}
	\left[
		\left|
		\big\|
		\frac{1}{\sqrt{k}}
		U\frac{\symbf{x}}{\left\|\symbf{x}\right\|}
		\big\|-1
		\right|\leq
		\varepsilon
		\right]
	\geq 1-
	2\exp
	\Big(-k\frac{\varepsilon^{2}}{16}\Big).
\end{equation}

\begin{definition}
	Sea $\left(\Omega,\Sigma,\mathbb{P}\right)$ un espacio de
	probabilidad y, para $k<d$, sea
	$U\colon\Omega\to\mathbb{R}^{k\times d}$ una matriz aleatoria con
	$U\sim\mathcal{N}\left(\symbf{0}_{k\times d},1\right)$.
	Además, sea $\omega\in\Omega$.
	Entonces, la aplicación
	\begin{equation}\label{eq:projection}
		\begin{aligned}
			T_{U\left(\omega\right)}\colon\mathbb{R}^{d} & \longrightarrow
			\mathbb{R}^{k}                                                 \\
			\symbf{x}                                    & \longmapsto
			\frac{1}{\sqrt{k}}
			U\left(\omega\right)\symbf{x},
		\end{aligned}
	\end{equation}
	se denomina \emph{proyección de Johnson-Lindenstrauss}.
\end{definition}
Estrictamente, la transformación $T_{U\left(\omega\right)}$ no es una
proyección en el sentido del álgebra lineal, pues no necesariamente
son matrices ortonormales.
Con esta notación, se puede reformular la conclusión del
Teorema~\ref{thm:projection} como sigue:
\begin{math}
	\forall\symbf{x}\in\mathbb{R}^{d}\setminus
	\left\{\symbf{0}_{d}\right\}
\end{math}
con $0<\varepsilon\leq 1$ y
$U\sim\mathcal{N}\left(\symbf{0}_{k\times d},1\right)$:
\begin{equation}
	\mathbb{P}
	\left[
		\left|
		\big\|T_{U\left(\omega\right)}\left(\symbf{x}\right)\big\|-
		\left\|\symbf{x}\right\|
		\right|\leq
		\varepsilon\left\|\symbf{x}\right\|
		\right]
	\geq 1-
	2\exp
	\Big(
	-k\frac{\varepsilon^{2}}{16}
	\Big),
\end{equation}
se observa que con alta probabilidad
\begin{equation*}
	\left\|
	T_{U\left(\omega\right)}\left(\symbf{x}\right)
	\right\|
	\approx
	\left\|x\right\|.
\end{equation*}
Por tanto, de cierta forma, la proyección de Johnson-Lindenstrauss
mantiene prácticamente invariante la norma de cada
$\symbf{x}\in\mathbb{R}^{d}$.

\section{El teorema de Johnson-Lindenstrauss}

El siguiente teorema de Johnson-Lindenstrauss en el contexto
probabilístico establece lo que sucede cuando consideramos
las distancias entre pares de puntos tras usar la proyección de
Johnson-Lindenstrauss.

\begin{theorem}[Teorema de Johnson-Lindenstrauss~\cite{Wegner2024}]\label{thm:johnson-lindenstrauss}
	Sea $\left(\Omega,\Sigma,\mathbb{P}\right)$ un espacio de
	probabilidad y sea además $0<\varepsilon<1$, $n\geq 1$ y
	$k\geq\frac{48}{\varepsilon^2}\log n$.
	Sea $U\colon\Omega\to\mathbb{R}^{k\times d}$ una matriz aleatoria
	con $U\sim\mathcal{N}\left(\symbf{0}_{k\times d},1\right)$.
	Entonces, para cada
	\begin{math}
		\left\{
		\symbf{x}_{1},
		\dotsc,
		\symbf{x}_{n}
		\right\}\subset\mathbb{R}^{d}
	\end{math}
	se cumple la siguiente estimación:
	\begin{equation*}
		\mathbb{P}
		\left[
			\forall i,j:
			\left(1-\varepsilon\right)
			\left\|
			\symbf{x}_{i}-
			\symbf{x}_{j}
			\right\|\leq
			\left\|
			T_{U\left(\omega\right)}\big(\symbf{x}_{i}\big)-
			T_{U\left(\omega\right)}\big(\symbf{x}_{j}\big)
			\right\|\leq
			\left(1+\varepsilon\right)
			\left\|
			\symbf{x}_{i}-
			\symbf{x}_{j}
			\right\|
			\right]\geq
		1-\frac{1}{n}.
	\end{equation*}
\end{theorem}

\begin{proof}
	Para cada $i\neq j$ se define
	$\symbf{x}\coloneqq\symbf{x}_i-\symbf{x}_{j}$ y se aplica el
	Teorema~\ref{thm:projection} para acotar $\mathbb{P}_{i,j}$ como
	sigue:
	\begin{align}
		\mathbb{P}_{i,j} & \coloneqq
		\mathbb{P}
		\left[
			\left(1-\varepsilon\right)
			\left\|
			\symbf{x}_{i}-
			\symbf{x}_{j}
			\right\|\leq
			\left\|
			T_{U\left(\omega\right)}\left(\symbf{x}_{i}\right)-
			T_{U\left(\omega\right)}\left(\symbf{x}_{j}\right)\right\|
			\leq
			\left(1+\varepsilon\right)
			\left\|
			\symbf{x}_{i}-
			\symbf{x}_{j}
			\right\|
		\right].        \notag                  \\
		\mathbb{P}_{i,j} & =
		\mathbb{P}
		\left[
			\left(1-\varepsilon\right)
			\left\|\symbf{x}\right\|\leq
			\left\|
			T_{U\left(\omega\right)}\left(\symbf{x}\right)
			\right\|\leq
			\left(1+\varepsilon\right)
			\left\|\symbf{x}\right\|
		\right].           \notag               \\
		\mathbb{P}_{i,j} & =
		\mathbb{P}
		\left[
			\left|
			\big\|
			T_{U\left(\omega\right)}
			\left(\symbf{x}\right)
			\big\|-
			\left\|
			\symbf{x}
			\right\|
			\right|\leq
		\varepsilon\|\symbf{x}\|\right]. \notag \\
		\mathbb{P}_{i,j} & \geq
		1-
		2\exp
		\Big(
		-k\frac{\varepsilon^{2}}{16}
		\Big).\label{eq:inequality}
	\end{align}
	Puesto que hay $n$ puntos en el conjunto
	$\left\{\symbf{x}_{1},\dotsc,\symbf{x}_{n}\right\}$, existen
	$\binom{n}{2}\leq\frac{n^{2}}{2}$ pares distintos de puntos, por
	tanto, se puede acotar la probabilidad $\mathbb{P}$ teniendo en
	cuenta~\eqref{eq:inequality} como sigue:
	\begin{align*}
		 &
		\mathbb{P}
		\left[
			\forall i,j:
			\left(1-\varepsilon\right)
			\left\|
			\symbf{x}_{i}-
			\symbf{x}_{j}
			\right\|\leq
			\left\|
			T_{U\left(\omega\right)}\left(\symbf{x}_{i}\right)-
			T_{U\left(\omega\right)}\left(\symbf{x}_{j}\right)
			\right\|\leq
			\left(1+\varepsilon\right)
			\left\|\symbf{x}_{i}-\symbf{x}_{j}\right\|
		\right]                                    \\
		 & =
		1-\mathbb{P}
		\left[
			\bigcup_{i,j}
			\left\{
			\left|
			\left\|
			T_{U\left(\omega\right)}\left(\symbf{x}_{i}\right)-
			T_{U\left(\omega\right)}\left(\symbf{x}_{j}\right)
			\right\|-
			\left\|
			\symbf{x}_{i}-
			\symbf{x}_{j}
			\right\|
			\right|>
			\varepsilon
			\left\|
			\symbf{x}_{i}-\symbf{x}_{j}
			\right\|
			\right\}
		\right].                                   \\
		 & \geq
		1-
		\sum_{i\neq j}
		\left(1-\mathbb{P}_{i,j}\right)\geq
		1-\sum_{i\neq j}
		2\exp\Big(-k\frac{\varepsilon^2}{16}\Big)\geq
		1-\frac{n^2}{2}
		2\exp\Big(-k\frac{\varepsilon^2}{16}\Big). \\
		 & \geq
		1-n^{2}
		\exp
		\left[
			-\Big(\frac{48}{\varepsilon^2}\log n\Big)
			\frac{\varepsilon^2}{16}
			\right]\geq
		1-
		n^{2}
		\exp
		\left[\log\left(n^{-3}\right)\right]=
		1-\frac{1}{n}.
	\end{align*}
\end{proof}

\begin{remark}
	Para aplicar el Teorema~\ref{thm:johnson-lindenstrauss}
	correctamente se debe interpretar la relación que existe entre las
	distintas constantes que aparecen en el enunciado:
	$n$ tamaño de la muestra en consideración, $d$ dimensión original
	del conjunto de puntos y $k$ dimensión a la que se pretende
	proyectar los puntos.
	En particular, $k\geq\frac{48}{\varepsilon^{2}}\log n$ crece de
	forma logarítmica con respecto a $n$ e inversamente proporcional al
	cuadrado de $\varepsilon$, lo cual quiere decir que conforme más
	grande sea el tamaño muestral $n$ y más pequeño sea $\varepsilon$
	más grande será la dimensión $k$ a la que se debe proyectar.
\end{remark}

\begin{remark}
	La clave está en que la dimensión proyectada $k$ no depende de la
	dimensión original $d$, solo depende de $n$ y $\varepsilon$.
	El Teorema~\ref{thm:johnson-lindenstrauss} garantiza que, al
	proyectar un conjunto de $n$ puntos desde un espacio de alta
	dimensión a un subespacio euclídeo de dimensión
	$k\geq\frac{48}{\varepsilon^{2}}\log n$, es posible conservar las
	distancias entre todos los pares de puntos con un error de a lo
	sumo $\varepsilon$, y esto implica que la estructura geométrica del
	conjunto se preserva casi idénticamente por la proyección de
	Johnson-Lindenstrauss $T_{U\left(\omega\right)}$.
\end{remark}

\begin{remark}
	El Teorema~\ref{thm:johnson-lindenstrauss} proporciona resultados
	relevantes en el caso de que $d$ sea de una dimensión moderadamente
	grande en comparación con $k$ que es del orden de $\log n$.
	Esto nos permite incluso considerar conjuntos de puntos más grandes
	que la dimensión del espacio $d$ que se podrán proyectar a un
	espacio de dimensión menor sin perder su geometría.
\end{remark}

\begin{remark}
	Además, los resultados del Teorema~\ref{thm:johnson-lindenstrauss}
	son totalmente independientes del conjunto original de puntos
	\begin{math}
		\left\{\symbf{x}_{1},\dotsc,\symbf{x}_{n}\right\}\subset
		\mathbb{R}^{d}
	\end{math}
	o de la distribución que los generó (si la tuviesen).
	De hecho, únicamente son dependientes del tamaño del conjunto sobre
	el que se quieran aplicar.
\end{remark}
